\chapter{损耗与增益}
\label{ch:attenuation}


光在玻璃介质中传播时，通常会因为杂质或晶格缺陷把部分电磁场散射到原本传播方向之外，从而自然地产生损耗。由于光强降低会削弱非线性光学效应，理解损耗的作用是理解式~\ref{eq:GNLSE} 的重要前提。


\section{损耗的物理来源}
当光在光纤这类波导中传播时，损耗可能来自玻璃中的化学杂质、晶格缺陷，或者波导中的急弯。此外，如果光的载波频率恰好与晶格中某种化学键的振动频率匹配，光能还可能直接转化为热。关于光如何在理想介质中传播的直观图像，可参考 \href{https://www.youtube.com/watch?v=QCX62YJCmGk&list=PLZHQObOWTQDMKqfyUvG2kTlYt-QQ2x-ui}{这个视频系列}。


\section{如何描述损耗}

考虑式~\ref{eq:GNLSE} 中除 $\alpha$ 外其余参数都为零的情形，此时有

\begin{align}
    \label{eq:attenuation}
    \partial_z\A &= \frac{\alpha}{2} \A \\ \nonumber
    \A(z,T)&=\A(0,T)\exp\left( \frac{\alpha}{2}z \right).
\end{align}
当 $\alpha<0$ 时，式~\ref{eq:attenuation} 表明场随传播距离衰减；当 $\alpha>0$ 时，则表示场经历增益。注意场的功率与场幅值的绝对值平方成正比，因此
\begin{align}
    \label{eq:attenuation_power}
    P(z,t)=|\A(z,t)|^2&=|\A(0,t)|^2 \exp\left( \alpha z \right).
\end{align}
在式~\ref{eq:GNLSE}、式~\ref{eq:attenuation} 和式~\ref{eq:attenuation_power} 中，$\alpha$ 表示“功率损耗系数”，也就是说它决定了场的\emph{功率}在传播一定距离后改变了多少个 $e$ 倍。也有一些作者把 $\alpha$ 定义为“场损耗系数”，用来描述\emph{场本身}变化了多少个 $e$ 倍；在那种定义下，式~\ref{eq:attenuation_power} 指数中的项应写成 $2\alpha z$，而不是 $\alpha z$。

\section{dB 标度}
\begin{table}[]
    \centering
    \begin{tabular}{ c|c|c|c|c|c|c|c|c|c|c|c|c|c|c|c|c|c }
 Scaling factor &0.01&0.05 & 0.1 &0.125 &0.25&0.4 &0.5 & 0.8&1&1.25 &2 &2.5 &4 &8 &10 &20 & \\  \hline
 dB change & -20 &-13 &-10 & -9&  -6&-4& -3&-1&0&1 &3 &4 &6 &9 & 10&13
\end{tabular}
    \caption{若干缩放因子及其对应的 dB 变化。}
    \label{tab:dB}
\end{table}
在实际应用里，损耗和增益通常用“分贝”（dB）来表示。比如，一个信号起初功率为 60~mW，穿过某介质后只剩 130~$\mu$W，那么它的功率变化为
\begin{align}
    \Delta P [dB] &= 10 \log_{10}\left(\frac{P_{final}}{P_{initial}} \right)= 10 \log_{10}\left(\frac{0.130 mW}{60 mW} \right) = -26.64 dB.
\end{align}
反过来，如果之后再把它放大 15~dB，那么新的功率为
\begin{align}
\label{eq:boost}
    P_{final}&=P_{initial}\cdot10^{\frac{\Delta dB}{10}}=0.13mW\cdot10^{\frac{15}{10}} = 4.11 mW.
\end{align}


注意 $10\log_{10}(10)=10$~dB，$10\log_{10}(0.1)=-10$~dB，$10\log_{10}(2)\approx3$~dB，而 $10\log_{10}(0.5)\approx-3$~dB。因此，如果功率被放大了 $20=10\cdot 2$ 倍，那么对应的 dB 增量就是 $10\log_{10}(10\cdot 2)=10\log_{10}(10)+10\log_{10}(2)=13$~dB。更多例子见表~\ref{tab:dB}。

损耗系数 $\alpha$ 常以 dB/km 为单位给出。例如，通信单模光纤在约 193~THz（约 1550~nm）附近的典型损耗约为 $-0.22$~dB/km。这意味着一段 100~km 的光纤会使光功率降低 $-22$~dB，也就是缩小约 158.5 倍。如果 $\alpha=-0.22$~dB/km，想把它换算为式~\ref{eq:GNLSE} 中需要用到的“每公里多少个 $e$ 倍”的形式（假设 $z$ 以 km 为单位），则有
\begin{align}
    \exp\left(\alpha_{\text{Factors of $e$ per km}}z\right) &= 10^{ \frac{\alpha_{dB/km}}{10}z} \\ \nonumber
    \alpha_{\text{Factors of $e$ per km}}z &= \ln\left(10^{ \frac{\alpha_{dB/km}}{10}z} \right)\\ \nonumber
    &= \frac{\alpha_{dB/km}}{10}z \ln(10)\\ \nonumber
    \alpha_{\text{Factors of $e$ per km}} &= 0.23\cdot \alpha_{dB/km}\\ \nonumber
    \alpha_{\text{Factors of $e$ per km}}\cdot4.343&=\alpha_{dB/km}.
\end{align}


\section{用 dBm 表示功率}
光信号功率还常用 “dBm” 表示，即“相对于 1~mW 的分贝值”。例如，40~mW 对应
\begin{align}
    P [dBm] &= 10\cdot\log_{10}\left(\frac{40mW}{1mW} \right)=16dBm,
\end{align}
反之，
\begin{align}
    \label{eq:dBm_rev}
    P [mW] &= 10^{\frac{16dBm}{10}}mW=40mW.
\end{align}


使用 dBm 很方便，因为即使是可探测的弱信号，功率也可能只有几 nW；而经过例如 \href{https://youtu.be/Eh5CHRWFT-M}{啁啾脉冲放大} 之后的脉冲峰值功率则可能达到 MW 甚至 GW~\cite{Chirp_STRICKLAND_Nobel_prize}。相比直接使用式~\ref{eq:boost}，dBm 还让增益与损耗的累计更加直接。例如，一个初始功率为 $0.13$~mW，也就是 $-8.86$~dBm 的信号，若被放大 15~dB，则末端功率就是 $-8.86$~dBm$+15$~dB$=6.13$~dBm。


\section{常见的 dB 与 dBm 错误}

\begin{enumerate}
    \item \textbf{把 dBm 和 dBm 直接相加}：Alice 想求两台激光器总功率的 dBm 值。她测得第一台是 3~dBm，第二台是 6~dBm。
    \begin{itemize}
    \item[\redcross] Alice 直接算 $P_{tot} [dBm] = 3dBm + 6dBm = 9dBm$。这是错的，因为 3~dBm = 2~mW，6~dBm = 4~mW，但 9~dBm = 8~mW；直接把 dBm 相加，等价于在线性单位下把功率相乘：$2~mW\cdot4~mW = 8~mW^2$！
    \item[\greencheck] Alice 先把 3~dBm 和 6~dBm 转成线性单位（分别为 2~mW 和 4~mW），相加得到总功率 6~mW，再换回 dBm：$10\cdot\log_{10}(6mW/1mW)=7.78$~dBm。
    \end{itemize}

    \item \textbf{混淆 dB 和 dBm}：Bob 想知道某台激光器的功率是多少 mW。他的一位同事告诉他测得功率是 “13~dB”。
    \begin{itemize}
    \item[\redcross] Bob 把 $13~dB$ 当成式~\ref{eq:dBm_rev} 中的 dBm 值来代入，算出 20~mW。错在“13~dB”是无量纲比值，并不是一个绝对功率值。
    \item[\greencheck] Bob 会先追问同事是不是其实想说的是 “13~dBm”。对方可能只是口误，也可能真的是功率计被设置成相对于某个参考功率显示 “13~dB”。
    \end{itemize}

   \item \textbf{再次混淆 dB 和 dBm}：Charlie 测得某激光器功率为 10~dBm。随后同事让他把功率“降低 3~dBm”。
    \begin{itemize}
    \item[\redcross] Charlie 直接把输出从 10~dBm 调到 7~dBm。这样做实际上是把功率降低了 3~dB，而“降低到 3~dBm”则意味着功率应变成 2~mW。
    \item[\greencheck] Charlie 应该先确认对方是想说降低 3~dB，还是确实想把输出调到 2~mW。前者显然更常见，因为后者虽然严格说并非错误，但表达上很不规范，也容易引起误解。
    \end{itemize}

    \item \textbf{错误求平均}：David 想求三台激光器功率的平均值，三者分别为 $-1$~dBm、4~dBm 和 6~dBm。
    \begin{itemize}
    \item[\redcross] David 直接算 $P_{avg}=(-1~dBm+4~dBm+6~dBm)/3=4dBm$。这得到的不是通常意义下的算术平均，而是这三个值的 \href{https://en.wikipedia.org/wiki/Geometric_mean}{几何平均}：$(-1~dBm+4~dBm+6~dBm)/3 = (0.79~mW\cdot2.51~mW\cdot7.94~mW)^{1/3}=2.5~mW = 4dBm$。这个量在某些场景下也有意义，但并不是 David 现在要的“普通平均值”。比如在一个团队共同表征某个待售光学产品性能时，平均功率到底是按哪种方式算，必须统一清楚。
    \item[\greencheck] David 应先把三个功率都转换成线性单位，再求平均：$(0.79~mW+2.51~mW+7.94~mW)/3=3.75~mW=5.74~dBm$。
    \end{itemize}



\end{enumerate}
